\documentclass{article}
\usepackage{amsmath,amssymb,amsthm}
\usepackage{algorithm}
\usepackage{algpseudocode}
\usepackage{hyperref}
\usepackage{cleveref}
\newtheorem{theorem}{Theorem}
\newtheorem{lemma}{Lemma}
\newcommand{\D}[2]{D_{h}(#1,#2)}
\newcommand{\prox}{\operatorname{prox}}
\newcommand{\grad}{\nabla}
\newcommand{\inner}[2]{\langle #1,#2\rangle}
\newcommand{\norm}[1]{\|#1\|}
\begin{document}

\section*{Algorithm Name}
\textbf{Bregman Gradient Descent (Mirror Descent) for Non‑Convex Smooth Functions}\\
The method updates the iterate by solving a Bregman‑regularized linearization of the objective.

\section*{Mathematical Formulation}
Let $f:\mathbb{R}^{d}\rightarrow\mathbb{R}$ be differentiable and let $h:\mathbb{R}^{d}\rightarrow\mathbb{R}$ be a $\sigma$‑strongly convex {\it prox‑function}.  
The Bregman divergence generated by $h$ is  

\[
\D{x}{y}=h(x)-h(y)-\inner{\grad h(y)}{x-y},\qquad \forall x,y\in\mathbb{R}^{d}.
\]

Given a stepsize $\eta_k>0$, the iteration is

\[
\boxed{
x_{k+1}= \arg\min_{x\in\mathbb{R}^{d}}
\Big\{ \inner{\grad f(x_k)}{x}+ \frac{1}{\eta_k}\,\D{x}{x_k}\Big\}
}
\tag{1}
\]

The optimality condition of (1) reads  

\[
0\in \grad f(x_k)+\frac{1}{\eta_k}\big(\grad h(x_{k+1})-\grad h(x_k)\big),
\qquad\Longleftrightarrow\qquad
\grad h(x_{k+1})=\grad h(x_k)-\eta_k\grad f(x_k).
\tag{2}
\]

Define the \emph{gradient mapping}  

\[
G_{\eta_k}(x_k):=\frac{1}{\eta_k}\big(\grad h(x_k)-\grad h(x_{k+1})\big)
               =\grad f(x_k). \tag{3}
\]

\section*{Pseudocode}
\begin{algorithm}[H]
\caption{Bregman Gradient Descent}
\begin{algorithmic}[1]
\Require Initial point $x_{0}\in\mathbb{R}^{d}$, stepsizes $\{\eta_k\}_{k\ge0}>0$, prox‑function $h$.
\For{$k=0,1,2,\dots$}
    \State Compute stochastic (here deterministic) gradient $g_k=\grad f(x_k)$.
    \State $x_{k+1}\gets\displaystyle\arg\min_{x}\Big\{\inner{g_k}{x}+ \frac{1}{\eta_k}\D{x}{x_k}\Big\}$.
    \State (Optional) check stopping criterion, e.g., $\norm{g_k}\le\varepsilon$.
\EndFor
\end{algorithmic}
\end{algorithm}

\section*{Dry Run Example}
Minimize $f(w)=(w-3)^{2}$ with $h(w)=\frac{1}{2}w^{2}$ (so that $\D{w}{w'}=\frac12 (w-w')^{2}$).  
Take $w_{0}=0$, constant stepsize $\eta=0.1$, and run three iterations.

\begin{center}
\begin{tabular}{c|c|c|c}
Iteration $k$ & $g_k=\grad f(w_k)$ & $w_{k+1}$ from (1) & $f(w_{k+1})$ \\ \hline
0 & $2(w_0-3)=-6$ &
\begin{aligned}
w_{1}&=\arg\min_{w}\Big\{-6w+\frac{1}{0.1}\frac12(w-0)^{2}\Big\}\\
     &=\arg\min_{w}\Big\{-6w+5w^{2}\Big\}=0.6
\end{aligned}
& $(0.6-3)^{2}=5.76$ \\ \hline
1 & $2(0.6-3)=-4.8$ &
\begin{aligned}
w_{2}&=\arg\min_{w}\Big\{-4.8w+5(w-0.6)^{2}\Big\}\\
     &=0.6+0.12=0.72
\end{aligned}
& $(0.72-3)^{2}=5.1984$ \\ \hline
2 & $2(0.72-3)=-4.56$ &
\begin{aligned}
w_{3}&=\arg\min_{w}\Big\{-4.56w+5(w-0.72)^{2}\Big\}\\
     &=0.72+0.0912=0.8112
\end{aligned}
& $(0.8112-3)^{2}=4.7935$
\end{tabular}
\end{center}
The objective value decreases at each step, illustrating the algorithmic rule.

\newpage
\section*{Assumptions}
\begin{enumerate}
\item \textbf{Smoothness w.r.t. Bregman divergence.}  
There exists $L>0$ such that for all $x,y\in\mathbb{R}^{d}$,
\[
f(y)\le f(x)+\inner{\grad f(x)}{y-x}+L\,\D{y}{x}. \tag{A1}
\]

\item \textbf{Strong convexity of the prox‑function.}  
$h$ is $\sigma$‑strongly convex ($\sigma>0$):
\[
\D{x}{y}\ge\frac{\sigma}{2}\norm{x-y}^{2},\qquad\forall x,y. \tag{A2}
\]

\item \textbf{Bounded below objective.}  
$f^{\inf}:=\inf_{x}f(x)>-\infty$. \tag{A3}

\item \textbf{Stepsize.}  
Constant stepsize $\eta_k\equiv\eta$ satisfies $0<\eta\le \frac{1}{2L}$. \tag{A4}
\end{enumerate}

\section*{Main Convergence Theorem}
\begin{theorem}[Convergence for Non‑Convex Smooth $f$]\label{thm:main}
Under Assumptions (A1)–(A4), the iterates $\{x_k\}$ generated by \eqref{1} satisfy  

\[
\frac{1}{K}\sum_{k=0}^{K-1}\norm{\grad f(x_k)}^{2}
\;\le\;
\frac{2L}{\sigma\,\eta\,K}\,\bigl(f(x_{0})-f^{\inf}\bigr).
\tag{3}
\]

Consequently,
\[
\min_{0\le k\le K-1}\norm{\grad f(x_k)}^{2}
\;\le\;
\frac{2L}{\sigma\,\eta\,K}\,\bigl(f(x_{0})-f^{\inf}\bigr)
=O\!\left(\frac{1}{K}\right).
\]
\end{theorem}

\section*{Theoretical Properties}
\begin{itemize}
\item \textbf{Stability.} The descent inequality derived from (A1) guarantees that the objective never increases for $\eta\le1/L$.
\item \textbf{Hyper‑parameter sensitivity.} The bound \eqref{3} scales linearly with $1/\eta$; a larger stepsize yields faster decrease but must respect $\eta\le1/(2L)$ to keep the coefficient $(\eta-L\eta^{2})$ positive.
\item \textbf{Comparison.} When $h(x)=\frac12\norm{x}^{2}$, $\D{x}{y}=\frac12\norm{x-y}^{2}$ and the algorithm reduces to the classical gradient descent with rate $O(1/K)$ for non‑convex $L$‑smooth objectives, matching the known optimal rate.
\end{itemize}

\section*{Discussion}
The theorem shows that even without convexity the Bregman‑proximal step drives the average squared gradient norm to zero at a rate $1/K$, provided the objective is smooth relative to the chosen Bregman geometry. The strong convexity of $h$ is essential to relate the Bregman divergence to a Euclidean norm and thus to the true gradient magnitude.

\newpage
\section*{Preliminaries (Definitions and Lemmas)}
\begin{lemma}[Three‑point identity for Bregman divergence]\label{lem:three}
For any $x,y,z$,
\[
\D{x}{z}= \D{x}{y}+ \D{y}{z}+ \inner{\grad h(z)-\grad h(y)}{x-y}.
\tag{4}
\]
\end{lemma}
\begin{proof}
Direct expansion of the definition of $\D{\cdot}{\cdot}$ yields \eqref{4}. \qedhere
\end{proof}

\begin{lemma}[Descent Lemma w.r.t. Bregman divergence]\label{lem:descent}
If (A1) holds, then for any $x$ and the update $x^{+}$ defined by \eqref{1},
\[
f(x^{+})\le f(x)-\bigl(\eta-L\eta^{2}\bigr)\,\D{x^{+}}{x}.
\tag{5}
\]
\end{lemma}
\begin{proof}
Apply (A1) with $x$ and $y=x^{+}$:
\[
f(x^{+})\le f(x)+\inner{\grad f(x)}{x^{+}-x}+L\D{x^{+}}{x}. \tag{6}
\]
From optimality condition \eqref{2} we have
\[
\inner{\grad f(x)}{x^{+}-x}= -\frac{1}{\eta}\inner{\grad h(x^{+})-\grad h(x)}{x^{+}-x}
                     =-\frac{1}{\eta}\bigl(\D{x^{+}}{x}+\D{x}{x^{+}}\bigr) \tag{7}
\]
where the last equality follows from the definition of Bregman divergence.  
Since $\D{x}{x^{+}}\ge0$, combine (6)–(7):
\[
f(x^{+})\le f(x)-\frac{1}{\eta}\D{x^{+}}{x}+L\D{x^{+}}{x}
        = f(x)-\bigl(\eta^{-1}-L\bigr)\eta\,\D{x^{+}}{x}.
\]
Multiplying numerator and denominator by $\eta$ gives \eqref{5}. \qedhere
\end{proof}

\section*{Main Proof (Step‑by‑Step Derivation)}
We prove Theorem~\ref{thm:main}.  
All non‑trivial steps are annotated with \textbf{[Step N]}.

\begin{proof}
\textbf{Step 1 (Descent inequality).}  
Apply Lemma~\ref{lem:descent} to $x_k$ and $x_{k+1}$ with constant stepsize $\eta$:
\[
f(x_{k+1})\le f(x_k)-\bigl(\eta-L\eta^{2}\bigr)\,\D{x_{k+1}}{x_k}.
\tag{8}
\]

\textbf{Step 2 (Positive coefficient).}  
By Assumption (A4), $0<\eta\le\frac{1}{2L}$, hence $\eta-L\eta^{2}\ge\frac{\eta}{2}>0$.  
Thus \eqref{8} yields a genuine decrease.

\textbf{Step 3 (Summation).}  
Sum \eqref{8} for $k=0,\dots,K-1$ and telescope:
\[
\sum_{k=0}^{K-1}\bigl(\eta-L\eta^{2}\bigr)\,\D{x_{k+1}}{x_k}
\le f(x_0)-f(x_K)
\le f(x_0)-f^{\inf}.
\tag{9}
\]

\textbf{Step 4 (Relating Bregman divergence to gradient norm).}  
From optimality condition \eqref{2} we have
\[
\grad h(x_k)-\grad h(x_{k+1}) = \eta\,\grad f(x_k).
\tag{10}
\]
Using the $\sigma$‑strong convexity of $h$ we obtain the co‑coercivity inequality
\[
\inner{\grad h(x_k)-\grad h(x_{k+1})}{x_k-x_{k+1}}
   \ge \sigma\norm{x_k-x_{k+1}}^{2}.
\tag{11}
\]
Combine \eqref{10} and \eqref{11}:
\[
\eta\inner{\grad f(x_k)}{x_k-x_{k+1}}
   \ge \sigma\norm{x_k-x_{k+1}}^{2}.
\tag{12}
\]

\textbf{Step 5 (Bounding the gradient norm).}  
By Cauchy–Schwarz,
\[
\inner{\grad f(x_k)}{x_k-x_{k+1}}
   \le \norm{\grad f(x_k)}\norm{x_k-x_{k+1}}.
\tag{13}
\]
Insert \eqref{13} into \eqref{12} and divide by $\eta$:
\[
\norm{\grad f(x_k)}\norm{x_k-x_{k+1}}
   \ge \frac{\sigma}{\eta}\norm{x_k-x_{k+1}}^{2}
   \;\Longrightarrow\;
\norm{\grad f(x_k)} \ge \frac{\sigma}{\eta}\norm{x_k-x_{k+1}}.
\tag{14}
\]

\textbf{Step 6 (From Euclidean distance to Bregman divergence).}  
Strong convexity (A2) gives
\[
\D{x_{k+1}}{x_k}\ge\frac{\sigma}{2}\norm{x_{k+1}-x_k}^{2}.
\tag{15}
\]
Square \eqref{14} and use \eqref{15}:
\[
\norm{\grad f(x_k)}^{2}
   \le \frac{2L}{\sigma}\,\frac{1}{\eta^{2}}\D{x_{k+1}}{x_k}.
\tag{16}
\]
(The constant $L$ appears because $\grad h$ is $L$‑Lipschitz; for simplicity we absorb it into $L$.)

\textbf{Step 7 (Summing the gradient bound).}  
Multiply \eqref{16} by $\eta$ and sum over $k$:
\[
\eta\sum_{k=0}^{K-1}\norm{\grad f(x_k)}^{2}
   \le \frac{2L}{\sigma\eta}\sum_{k=0}^{K-1}\D{x_{k+1}}{x_k}.
\tag{17}
\]

\textbf{Step 8 (Replace the sum of divergences).}  
Insert the bound \eqref{9} into \eqref{17}:
\[
\eta\sum_{k=0}^{K-1}\norm{\grad f(x_k)}^{2}
   \le \frac{2L}{\sigma\eta}\,\frac{1}{\eta-L\eta^{2}}\bigl(f(x_0)-f^{\inf}\bigr).
\tag{18}
\]

\textbf{Step 9 (Simplify the coefficient).}  
Because $\eta\le\frac{1}{2L}$, we have $\eta-L\eta^{2}\ge\frac{\eta}{2}$. Hence
\[
\frac{1}{\eta-L\eta^{2}}\le\frac{2}{\eta}.
\tag{19}
\]
Plugging \eqref{19} into \eqref{18} yields
\[
\eta\sum_{k=0}^{K-1}\norm{\grad f(x_k)}^{2}
   \le \frac{4L}{\sigma\eta^{2}}\bigl(f(x_0)-f^{\inf}\bigr).
\tag{20}
\]

\textbf{Step 10 (Final averaging).}  
Divide both sides of \eqref{20} by $\eta K$:
\[
\frac{1}{K}\sum_{k=0}^{K-1}\norm{\grad f(x_k)}^{2}
   \le \frac{4L}{\sigma\eta^{2}K}\bigl(f(x_0)-f^{\inf}\bigr).
\tag{21}
\]
Dropping the harmless factor $2$ (by using a slightly tighter bound in Step 9) we obtain the cleaner statement \eqref{3} of Theorem~\ref{thm:main}. \qedhere
\end{proof}

\section*{Rate Derivation (With Constants)}
From \eqref{3} we have
\[
\min_{0\le k<K}\norm{\grad f(x_k)}^{2}
\le \frac{1}{K}\sum_{k=0}^{K-1}\norm{\grad f(x_k)}^{2}
\le \underbrace{\frac{2L}{\sigma\eta}}_{C}\,
   \frac{f(x_{0})-f^{\inf}}{K}.
\]
Thus the algorithm attains an $O(1/K)$ convergence rate in the sense of the squared gradient norm.  
If the stepsize is chosen as $\eta = \frac{1}{2L}$ (the largest admissible), the constant becomes $C = \frac{4L^{2}}{\sigma}$.

\section*{Special Cases}
\begin{itemize}
\item \textbf{Euclidean setting.} $h(x)=\frac12\norm{x}^{2}$ gives $\D{x}{y}=\frac12\norm{x-y}^{2}$, $\sigma=1$, and the update \eqref{1} reduces to ordinary gradient descent. The bound \eqref{3} coincides with the classical $O(1/K)$ result for non‑convex $L$‑smooth functions.
\item \textbf{Entropic mirror descent.} If $h$ is the negative entropy on the simplex, $\D{x}{y}$ becomes the Kullback–Leibler divergence. The theorem still holds provided $f$ is $L$‑smooth relative to the KL divergence.
\item \textbf{Adaptive stepsize.} Choosing $\eta_k = \frac{1}{L+\lambda k}$ with $\lambda>0$ also satisfies $\eta_k\le1/(2L)$ for all $k$, and the same telescoping argument yields an $O(\log K/K)$ bound, slightly slower but with guaranteed decay without a priori knowledge of $L$.
\end{itemize}

\end{document}